%%%%% PROJECT-SPECIFIC MACROS %%%%%
% Define your project-specific terms here
% Import with: %%%%% PROJECT-SPECIFIC MACROS %%%%%
% Define your project-specific terms here
% Import with: %%%%% PROJECT-SPECIFIC MACROS %%%%%
% Define your project-specific terms here
% Import with: %%%%% PROJECT-SPECIFIC MACROS %%%%%
% Define your project-specific terms here
% Import with: \input{commands/macros}
%
% INSTRUCTIONS:
% 1. Use \textsc{} for method/dataset/model names (small caps)
% 2. Add \xspace at the end for automatic spacing
% 3. Group related terms together
% 4. Use consistent naming conventions

\usepackage{xspace}

%%%%%%%%%%%%%%%%%%%%%%%%%%%%%%%%%%%%%%%%%%%%%%%%%%%%%%%%%%%%%%%%%
% METHOD NAMES
%%%%%%%%%%%%%%%%%%%%%%%%%%%%%%%%%%%%%%%%%%%%%%%%%%%%%%%%%%%%%%%%%

% Your method name (replace with actual name)
% Example: \newcommand{\ours}{{\bf YourMethod}\xspace}
% Example: \newcommand{\methodname}{\textsc{MethodName}\xspace}

%%%%%%%%%%%%%%%%%%%%%%%%%%%%%%%%%%%%%%%%%%%%%%%%%%%%%%%%%%%%%%%%%
% DATASET NAMES
%%%%%%%%%%%%%%%%%%%%%%%%%%%%%%%%%%%%%%%%%%%%%%%%%%%%%%%%%%%%%%%%%

% Dataset names (use \textsc for consistency)
% Example: \newcommand{\datasetone}{\textsc{Dataset-One}\xspace}
% Example: \newcommand{\datasettwo}{\textsc{Dataset-Two}\xspace}

%%%%%%%%%%%%%%%%%%%%%%%%%%%%%%%%%%%%%%%%%%%%%%%%%%%%%%%%%%%%%%%%%
% MODEL NAMES
%%%%%%%%%%%%%%%%%%%%%%%%%%%%%%%%%%%%%%%%%%%%%%%%%%%%%%%%%%%%%%%%%

% LLM and model names
% Example: \newcommand{\gpt}{\textsc{GPT-4o-mini}\xspace}
% Example: \newcommand{\llama}{\textsc{Llama-3.1-70B-Instruct}\xspace}
% Example: \newcommand{\claude}{\textsc{Claude-3.5-Sonnet}\xspace}

%%%%%%%%%%%%%%%%%%%%%%%%%%%%%%%%%%%%%%%%%%%%%%%%%%%%%%%%%%%%%%%%%
% BASELINE NAMES
%%%%%%%%%%%%%%%%%%%%%%%%%%%%%%%%%%%%%%%%%%%%%%%%%%%%%%%%%%%%%%%%%

% Baseline method names
% Example: \newcommand{\baseline}{\textsc{BaselineMethod}\xspace}
% Example: \newcommand{\fewshot}{\textsc{Few-shot}\xspace}

%%%%%%%%%%%%%%%%%%%%%%%%%%%%%%%%%%%%%%%%%%%%%%%%%%%%%%%%%%%%%%%%%
% ABBREVIATIONS
%%%%%%%%%%%%%%%%%%%%%%%%%%%%%%%%%%%%%%%%%%%%%%%%%%%%%%%%%%%%%%%%%

% Common abbreviations used in your paper
% Example: \newcommand{\ie}{i.e.,\xspace}
% Example: \newcommand{\eg}{e.g.,\xspace}
% Example: \newcommand{\cf}{cf.\xspace}
% Example: \newcommand{\wrt}{w.r.t.\xspace}

%%%%%%%%%%%%%%%%%%%%%%%%%%%%%%%%%%%%%%%%%%%%%%%%%%%%%%%%%%%%%%%%%
% DOMAIN-SPECIFIC TERMS
%%%%%%%%%%%%%%%%%%%%%%%%%%%%%%%%%%%%%%%%%%%%%%%%%%%%%%%%%%%%%%%%%

% Terms specific to your research domain
% Keep consistent formatting throughout the paper

%%%%%%%%%%%%%%%%%%%%%%%%%%%%%%%%%%%%%%%%%%%%%%%%%%%%%%%%%%%%%%%%%
% USAGE EXAMPLES
%%%%%%%%%%%%%%%%%%%%%%%%%%%%%%%%%%%%%%%%%%%%%%%%%%%%%%%%%%%%%%%%%
%
% GOOD:
%   \newcommand{\hypogenic}{\textsc{HypoGeniC}\xspace}
%   Usage: "We propose \hypogenic to generate hypotheses."
%   Result: "We propose HypoGeniC to generate hypotheses."
%
% GOOD:
%   \newcommand{\ours}{{\bf HypoGeniC}\xspace}
%   Usage: "Our method, \ours, outperforms baselines."
%   Result: "Our method, HypoGeniC, outperforms baselines."
%
% BAD (no \xspace):
%   \newcommand{\hypogenic}{\textsc{HypoGeniC}}
%   Usage: "\hypogenic outperforms..."
%   Result: "HypoGeniCoutperforms..." (missing space!)
%
%%%%%%%%%%%%%%%%%%%%%%%%%%%%%%%%%%%%%%%%%%%%%%%%%%%%%%%%%%%%%%%%%

%%%%% PROJECT-SPECIFIC DEFINITIONS (this paper) %%%%%
\newcommand{\ie}{i.e.,\xspace}
\newcommand{\eg}{e.g.,\xspace}
\newcommand{\pp}{\,pp\xspace}

% Experiments
\newcommand{\expA}{\textsc{Exp-A}\xspace}
\newcommand{\expB}{\textsc{Exp-B}\xspace}
\newcommand{\expD}{\textsc{Exp-D}\xspace}
\newcommand{\expE}{\textsc{Exp-E}\xspace}
\newcommand{\expF}{\textsc{Exp-F}\xspace}

% Datasets
\newcommand{\alpaca}{\textsc{Alpaca-cleaned}\xspace}
\newcommand{\wikitext}{\textsc{WikiText-103}\xspace}

% Models
\newcommand{\qhalf}{\textsc{Qwen2.5-0.5B}\xspace}
\newcommand{\qonefive}{\textsc{Qwen2.5-1.5B}\xspace}
\newcommand{\qthree}{\textsc{Qwen2.5-3B}\xspace}
\newcommand{\qseven}{\textsc{Qwen2.5-7B}\xspace}
\newcommand{\qfourteen}{\textsc{Qwen2.5-14B}\xspace}
\newcommand{\lone}{\textsc{Llama-3.2-1B}\xspace}
\newcommand{\lthree}{\textsc{Llama-3.2-3B}\xspace}
\newcommand{\leight}{\textsc{Llama-3.1-8B}\xspace}
\newcommand{\qwenthree}{\textsc{Qwen3-8B}\xspace}
\def\appref#1{appendix~\ref{#1}}

%
% INSTRUCTIONS:
% 1. Use \textsc{} for method/dataset/model names (small caps)
% 2. Add \xspace at the end for automatic spacing
% 3. Group related terms together
% 4. Use consistent naming conventions

\usepackage{xspace}

%%%%%%%%%%%%%%%%%%%%%%%%%%%%%%%%%%%%%%%%%%%%%%%%%%%%%%%%%%%%%%%%%
% METHOD NAMES
%%%%%%%%%%%%%%%%%%%%%%%%%%%%%%%%%%%%%%%%%%%%%%%%%%%%%%%%%%%%%%%%%

% Your method name (replace with actual name)
% Example: \newcommand{\ours}{{\bf YourMethod}\xspace}
% Example: \newcommand{\methodname}{\textsc{MethodName}\xspace}

%%%%%%%%%%%%%%%%%%%%%%%%%%%%%%%%%%%%%%%%%%%%%%%%%%%%%%%%%%%%%%%%%
% DATASET NAMES
%%%%%%%%%%%%%%%%%%%%%%%%%%%%%%%%%%%%%%%%%%%%%%%%%%%%%%%%%%%%%%%%%

% Dataset names (use \textsc for consistency)
% Example: \newcommand{\datasetone}{\textsc{Dataset-One}\xspace}
% Example: \newcommand{\datasettwo}{\textsc{Dataset-Two}\xspace}

%%%%%%%%%%%%%%%%%%%%%%%%%%%%%%%%%%%%%%%%%%%%%%%%%%%%%%%%%%%%%%%%%
% MODEL NAMES
%%%%%%%%%%%%%%%%%%%%%%%%%%%%%%%%%%%%%%%%%%%%%%%%%%%%%%%%%%%%%%%%%

% LLM and model names
% Example: \newcommand{\gpt}{\textsc{GPT-4o-mini}\xspace}
% Example: \newcommand{\llama}{\textsc{Llama-3.1-70B-Instruct}\xspace}
% Example: \newcommand{\claude}{\textsc{Claude-3.5-Sonnet}\xspace}

%%%%%%%%%%%%%%%%%%%%%%%%%%%%%%%%%%%%%%%%%%%%%%%%%%%%%%%%%%%%%%%%%
% BASELINE NAMES
%%%%%%%%%%%%%%%%%%%%%%%%%%%%%%%%%%%%%%%%%%%%%%%%%%%%%%%%%%%%%%%%%

% Baseline method names
% Example: \newcommand{\baseline}{\textsc{BaselineMethod}\xspace}
% Example: \newcommand{\fewshot}{\textsc{Few-shot}\xspace}

%%%%%%%%%%%%%%%%%%%%%%%%%%%%%%%%%%%%%%%%%%%%%%%%%%%%%%%%%%%%%%%%%
% ABBREVIATIONS
%%%%%%%%%%%%%%%%%%%%%%%%%%%%%%%%%%%%%%%%%%%%%%%%%%%%%%%%%%%%%%%%%

% Common abbreviations used in your paper
% Example: \newcommand{\ie}{i.e.,\xspace}
% Example: \newcommand{\eg}{e.g.,\xspace}
% Example: \newcommand{\cf}{cf.\xspace}
% Example: \newcommand{\wrt}{w.r.t.\xspace}

%%%%%%%%%%%%%%%%%%%%%%%%%%%%%%%%%%%%%%%%%%%%%%%%%%%%%%%%%%%%%%%%%
% DOMAIN-SPECIFIC TERMS
%%%%%%%%%%%%%%%%%%%%%%%%%%%%%%%%%%%%%%%%%%%%%%%%%%%%%%%%%%%%%%%%%

% Terms specific to your research domain
% Keep consistent formatting throughout the paper

%%%%%%%%%%%%%%%%%%%%%%%%%%%%%%%%%%%%%%%%%%%%%%%%%%%%%%%%%%%%%%%%%
% USAGE EXAMPLES
%%%%%%%%%%%%%%%%%%%%%%%%%%%%%%%%%%%%%%%%%%%%%%%%%%%%%%%%%%%%%%%%%
%
% GOOD:
%   \newcommand{\hypogenic}{\textsc{HypoGeniC}\xspace}
%   Usage: "We propose \hypogenic to generate hypotheses."
%   Result: "We propose HypoGeniC to generate hypotheses."
%
% GOOD:
%   \newcommand{\ours}{{\bf HypoGeniC}\xspace}
%   Usage: "Our method, \ours, outperforms baselines."
%   Result: "Our method, HypoGeniC, outperforms baselines."
%
% BAD (no \xspace):
%   \newcommand{\hypogenic}{\textsc{HypoGeniC}}
%   Usage: "\hypogenic outperforms..."
%   Result: "HypoGeniCoutperforms..." (missing space!)
%
%%%%%%%%%%%%%%%%%%%%%%%%%%%%%%%%%%%%%%%%%%%%%%%%%%%%%%%%%%%%%%%%%

%%%%% PROJECT-SPECIFIC DEFINITIONS (this paper) %%%%%
\newcommand{\ie}{i.e.,\xspace}
\newcommand{\eg}{e.g.,\xspace}
\newcommand{\pp}{\,pp\xspace}

% Experiments
\newcommand{\expA}{\textsc{Exp-A}\xspace}
\newcommand{\expB}{\textsc{Exp-B}\xspace}
\newcommand{\expD}{\textsc{Exp-D}\xspace}
\newcommand{\expE}{\textsc{Exp-E}\xspace}
\newcommand{\expF}{\textsc{Exp-F}\xspace}

% Datasets
\newcommand{\alpaca}{\textsc{Alpaca-cleaned}\xspace}
\newcommand{\wikitext}{\textsc{WikiText-103}\xspace}

% Models
\newcommand{\qhalf}{\textsc{Qwen2.5-0.5B}\xspace}
\newcommand{\qonefive}{\textsc{Qwen2.5-1.5B}\xspace}
\newcommand{\qthree}{\textsc{Qwen2.5-3B}\xspace}
\newcommand{\qseven}{\textsc{Qwen2.5-7B}\xspace}
\newcommand{\qfourteen}{\textsc{Qwen2.5-14B}\xspace}
\newcommand{\lone}{\textsc{Llama-3.2-1B}\xspace}
\newcommand{\lthree}{\textsc{Llama-3.2-3B}\xspace}
\newcommand{\leight}{\textsc{Llama-3.1-8B}\xspace}
\newcommand{\qwenthree}{\textsc{Qwen3-8B}\xspace}
\def\appref#1{appendix~\ref{#1}}

%
% INSTRUCTIONS:
% 1. Use \textsc{} for method/dataset/model names (small caps)
% 2. Add \xspace at the end for automatic spacing
% 3. Group related terms together
% 4. Use consistent naming conventions

\usepackage{xspace}

%%%%%%%%%%%%%%%%%%%%%%%%%%%%%%%%%%%%%%%%%%%%%%%%%%%%%%%%%%%%%%%%%
% METHOD NAMES
%%%%%%%%%%%%%%%%%%%%%%%%%%%%%%%%%%%%%%%%%%%%%%%%%%%%%%%%%%%%%%%%%

% Your method name (replace with actual name)
% Example: \newcommand{\ours}{{\bf YourMethod}\xspace}
% Example: \newcommand{\methodname}{\textsc{MethodName}\xspace}

%%%%%%%%%%%%%%%%%%%%%%%%%%%%%%%%%%%%%%%%%%%%%%%%%%%%%%%%%%%%%%%%%
% DATASET NAMES
%%%%%%%%%%%%%%%%%%%%%%%%%%%%%%%%%%%%%%%%%%%%%%%%%%%%%%%%%%%%%%%%%

% Dataset names (use \textsc for consistency)
% Example: \newcommand{\datasetone}{\textsc{Dataset-One}\xspace}
% Example: \newcommand{\datasettwo}{\textsc{Dataset-Two}\xspace}

%%%%%%%%%%%%%%%%%%%%%%%%%%%%%%%%%%%%%%%%%%%%%%%%%%%%%%%%%%%%%%%%%
% MODEL NAMES
%%%%%%%%%%%%%%%%%%%%%%%%%%%%%%%%%%%%%%%%%%%%%%%%%%%%%%%%%%%%%%%%%

% LLM and model names
% Example: \newcommand{\gpt}{\textsc{GPT-4o-mini}\xspace}
% Example: \newcommand{\llama}{\textsc{Llama-3.1-70B-Instruct}\xspace}
% Example: \newcommand{\claude}{\textsc{Claude-3.5-Sonnet}\xspace}

%%%%%%%%%%%%%%%%%%%%%%%%%%%%%%%%%%%%%%%%%%%%%%%%%%%%%%%%%%%%%%%%%
% BASELINE NAMES
%%%%%%%%%%%%%%%%%%%%%%%%%%%%%%%%%%%%%%%%%%%%%%%%%%%%%%%%%%%%%%%%%

% Baseline method names
% Example: \newcommand{\baseline}{\textsc{BaselineMethod}\xspace}
% Example: \newcommand{\fewshot}{\textsc{Few-shot}\xspace}

%%%%%%%%%%%%%%%%%%%%%%%%%%%%%%%%%%%%%%%%%%%%%%%%%%%%%%%%%%%%%%%%%
% ABBREVIATIONS
%%%%%%%%%%%%%%%%%%%%%%%%%%%%%%%%%%%%%%%%%%%%%%%%%%%%%%%%%%%%%%%%%

% Common abbreviations used in your paper
% Example: \newcommand{\ie}{i.e.,\xspace}
% Example: \newcommand{\eg}{e.g.,\xspace}
% Example: \newcommand{\cf}{cf.\xspace}
% Example: \newcommand{\wrt}{w.r.t.\xspace}

%%%%%%%%%%%%%%%%%%%%%%%%%%%%%%%%%%%%%%%%%%%%%%%%%%%%%%%%%%%%%%%%%
% DOMAIN-SPECIFIC TERMS
%%%%%%%%%%%%%%%%%%%%%%%%%%%%%%%%%%%%%%%%%%%%%%%%%%%%%%%%%%%%%%%%%

% Terms specific to your research domain
% Keep consistent formatting throughout the paper

%%%%%%%%%%%%%%%%%%%%%%%%%%%%%%%%%%%%%%%%%%%%%%%%%%%%%%%%%%%%%%%%%
% USAGE EXAMPLES
%%%%%%%%%%%%%%%%%%%%%%%%%%%%%%%%%%%%%%%%%%%%%%%%%%%%%%%%%%%%%%%%%
%
% GOOD:
%   \newcommand{\hypogenic}{\textsc{HypoGeniC}\xspace}
%   Usage: "We propose \hypogenic to generate hypotheses."
%   Result: "We propose HypoGeniC to generate hypotheses."
%
% GOOD:
%   \newcommand{\ours}{{\bf HypoGeniC}\xspace}
%   Usage: "Our method, \ours, outperforms baselines."
%   Result: "Our method, HypoGeniC, outperforms baselines."
%
% BAD (no \xspace):
%   \newcommand{\hypogenic}{\textsc{HypoGeniC}}
%   Usage: "\hypogenic outperforms..."
%   Result: "HypoGeniCoutperforms..." (missing space!)
%
%%%%%%%%%%%%%%%%%%%%%%%%%%%%%%%%%%%%%%%%%%%%%%%%%%%%%%%%%%%%%%%%%

%%%%% PROJECT-SPECIFIC DEFINITIONS (this paper) %%%%%
\newcommand{\ie}{i.e.,\xspace}
\newcommand{\eg}{e.g.,\xspace}
\newcommand{\pp}{\,pp\xspace}

% Experiments
\newcommand{\expA}{\textsc{Exp-A}\xspace}
\newcommand{\expB}{\textsc{Exp-B}\xspace}
\newcommand{\expD}{\textsc{Exp-D}\xspace}
\newcommand{\expE}{\textsc{Exp-E}\xspace}
\newcommand{\expF}{\textsc{Exp-F}\xspace}

% Datasets
\newcommand{\alpaca}{\textsc{Alpaca-cleaned}\xspace}
\newcommand{\wikitext}{\textsc{WikiText-103}\xspace}

% Models
\newcommand{\qhalf}{\textsc{Qwen2.5-0.5B}\xspace}
\newcommand{\qonefive}{\textsc{Qwen2.5-1.5B}\xspace}
\newcommand{\qthree}{\textsc{Qwen2.5-3B}\xspace}
\newcommand{\qseven}{\textsc{Qwen2.5-7B}\xspace}
\newcommand{\qfourteen}{\textsc{Qwen2.5-14B}\xspace}
\newcommand{\lone}{\textsc{Llama-3.2-1B}\xspace}
\newcommand{\lthree}{\textsc{Llama-3.2-3B}\xspace}
\newcommand{\leight}{\textsc{Llama-3.1-8B}\xspace}
\newcommand{\qwenthree}{\textsc{Qwen3-8B}\xspace}
\def\appref#1{appendix~\ref{#1}}

%
% INSTRUCTIONS:
% 1. Use \textsc{} for method/dataset/model names (small caps)
% 2. Add \xspace at the end for automatic spacing
% 3. Group related terms together
% 4. Use consistent naming conventions

\usepackage{xspace}

%%%%%%%%%%%%%%%%%%%%%%%%%%%%%%%%%%%%%%%%%%%%%%%%%%%%%%%%%%%%%%%%%
% METHOD NAMES
%%%%%%%%%%%%%%%%%%%%%%%%%%%%%%%%%%%%%%%%%%%%%%%%%%%%%%%%%%%%%%%%%

% Your method name (replace with actual name)
% Example: \newcommand{\ours}{{\bf YourMethod}\xspace}
% Example: \newcommand{\methodname}{\textsc{MethodName}\xspace}

%%%%%%%%%%%%%%%%%%%%%%%%%%%%%%%%%%%%%%%%%%%%%%%%%%%%%%%%%%%%%%%%%
% DATASET NAMES
%%%%%%%%%%%%%%%%%%%%%%%%%%%%%%%%%%%%%%%%%%%%%%%%%%%%%%%%%%%%%%%%%

% Dataset names (use \textsc for consistency)
% Example: \newcommand{\datasetone}{\textsc{Dataset-One}\xspace}
% Example: \newcommand{\datasettwo}{\textsc{Dataset-Two}\xspace}

%%%%%%%%%%%%%%%%%%%%%%%%%%%%%%%%%%%%%%%%%%%%%%%%%%%%%%%%%%%%%%%%%
% MODEL NAMES
%%%%%%%%%%%%%%%%%%%%%%%%%%%%%%%%%%%%%%%%%%%%%%%%%%%%%%%%%%%%%%%%%

% LLM and model names
% Example: \newcommand{\gpt}{\textsc{GPT-4o-mini}\xspace}
% Example: \newcommand{\llama}{\textsc{Llama-3.1-70B-Instruct}\xspace}
% Example: \newcommand{\claude}{\textsc{Claude-3.5-Sonnet}\xspace}

%%%%%%%%%%%%%%%%%%%%%%%%%%%%%%%%%%%%%%%%%%%%%%%%%%%%%%%%%%%%%%%%%
% BASELINE NAMES
%%%%%%%%%%%%%%%%%%%%%%%%%%%%%%%%%%%%%%%%%%%%%%%%%%%%%%%%%%%%%%%%%

% Baseline method names
% Example: \newcommand{\baseline}{\textsc{BaselineMethod}\xspace}
% Example: \newcommand{\fewshot}{\textsc{Few-shot}\xspace}

%%%%%%%%%%%%%%%%%%%%%%%%%%%%%%%%%%%%%%%%%%%%%%%%%%%%%%%%%%%%%%%%%
% ABBREVIATIONS
%%%%%%%%%%%%%%%%%%%%%%%%%%%%%%%%%%%%%%%%%%%%%%%%%%%%%%%%%%%%%%%%%

% Common abbreviations used in your paper
% Example: \newcommand{\ie}{i.e.,\xspace}
% Example: \newcommand{\eg}{e.g.,\xspace}
% Example: \newcommand{\cf}{cf.\xspace}
% Example: \newcommand{\wrt}{w.r.t.\xspace}

%%%%%%%%%%%%%%%%%%%%%%%%%%%%%%%%%%%%%%%%%%%%%%%%%%%%%%%%%%%%%%%%%
% DOMAIN-SPECIFIC TERMS
%%%%%%%%%%%%%%%%%%%%%%%%%%%%%%%%%%%%%%%%%%%%%%%%%%%%%%%%%%%%%%%%%

% Terms specific to your research domain
% Keep consistent formatting throughout the paper

%%%%%%%%%%%%%%%%%%%%%%%%%%%%%%%%%%%%%%%%%%%%%%%%%%%%%%%%%%%%%%%%%
% USAGE EXAMPLES
%%%%%%%%%%%%%%%%%%%%%%%%%%%%%%%%%%%%%%%%%%%%%%%%%%%%%%%%%%%%%%%%%
%
% GOOD:
%   \newcommand{\hypogenic}{\textsc{HypoGeniC}\xspace}
%   Usage: "We propose \hypogenic to generate hypotheses."
%   Result: "We propose HypoGeniC to generate hypotheses."
%
% GOOD:
%   \newcommand{\ours}{{\bf HypoGeniC}\xspace}
%   Usage: "Our method, \ours, outperforms baselines."
%   Result: "Our method, HypoGeniC, outperforms baselines."
%
% BAD (no \xspace):
%   \newcommand{\hypogenic}{\textsc{HypoGeniC}}
%   Usage: "\hypogenic outperforms..."
%   Result: "HypoGeniCoutperforms..." (missing space!)
%
%%%%%%%%%%%%%%%%%%%%%%%%%%%%%%%%%%%%%%%%%%%%%%%%%%%%%%%%%%%%%%%%%

%%%%% PROJECT-SPECIFIC DEFINITIONS (this paper) %%%%%
\newcommand{\ie}{i.e.,\xspace}
\newcommand{\eg}{e.g.,\xspace}
\newcommand{\pp}{\,pp\xspace}

% Experiments
\newcommand{\expA}{\textsc{Exp-A}\xspace}
\newcommand{\expB}{\textsc{Exp-B}\xspace}
\newcommand{\expD}{\textsc{Exp-D}\xspace}
\newcommand{\expE}{\textsc{Exp-E}\xspace}
\newcommand{\expF}{\textsc{Exp-F}\xspace}

% Datasets
\newcommand{\alpaca}{\textsc{Alpaca-cleaned}\xspace}
\newcommand{\wikitext}{\textsc{WikiText-103}\xspace}

% Models
\newcommand{\qhalf}{\textsc{Qwen2.5-0.5B}\xspace}
\newcommand{\qonefive}{\textsc{Qwen2.5-1.5B}\xspace}
\newcommand{\qthree}{\textsc{Qwen2.5-3B}\xspace}
\newcommand{\qseven}{\textsc{Qwen2.5-7B}\xspace}
\newcommand{\qfourteen}{\textsc{Qwen2.5-14B}\xspace}
\newcommand{\lone}{\textsc{Llama-3.2-1B}\xspace}
\newcommand{\lthree}{\textsc{Llama-3.2-3B}\xspace}
\newcommand{\leight}{\textsc{Llama-3.1-8B}\xspace}
\newcommand{\qwenthree}{\textsc{Qwen3-8B}\xspace}
\def\appref#1{appendix~\ref{#1}}
